
\begin{proposition}
    Let \(L\) be a finite-dimensional semisimple Lie algebra over an algebraically
    closed field \(F\) of characteristic \(0\). Let
    \[
        L=H\oplus \bigoplus_{\alpha\in\Delta} L_\alpha
    \]
    be the root space decomposition relative to a Cartan subalgebra \(H\).
    For each \(\alpha\in\Delta\), choose root vectors
    \[
        0\neq x_\alpha\in L_\alpha,\qquad
        0\neq y_\alpha\in L_{-\alpha},
    \]
    and set
    \[
        h_\alpha=[x_\alpha,y_\alpha],
        \qquad
        \alpha(h_\alpha)=2.
    \]
    Then
    \[
        \mathfrak{s}_\alpha
        :=
        F x_\alpha\oplus F h_\alpha\oplus F y_\alpha
        \cong \mathfrak{sl}_2(F).
    \]
    Moreover:
    \begin{enumerate}
        \item
              \[
                  F\alpha\cap\Delta=\{-\alpha,\alpha\},
                  \qquad
                  \dim L_\alpha=1.
              \]

        \item
              If \(\alpha,\beta\in\Delta\) and \(\beta\neq \pm\alpha\), then
              \[
                  \frac{2(\alpha,\beta)}{(\alpha,\alpha)}\in\mathbb{Z}.
              \]
              Furthermore, there exist \(p,q\in\mathbb{Z}_{\geq 0}\) such that the
              roots of the form \(\beta+i\alpha\) are precisely
              \[
                  \beta-q\alpha,\,
                  \beta-(q-1)\alpha,\,
                  \ldots,\,
                  \beta-\alpha,\,
                  \beta,\,
                  \beta+\alpha,\,
                  \ldots,\,
                  \beta+p\alpha,
              \]
              and
              \[
                  q-p
                  =
                  \frac{2(\alpha,\beta)}{(\alpha,\alpha)}.
              \]
              In particular, if \(\alpha+\beta\in\Delta\), then
              \[
                  [L_\alpha,L_\beta]=L_{\alpha+\beta}.
              \]
    \end{enumerate}
\end{proposition}

\begin{proof}
    We regard \(L\) as an \(\mathfrak{s}_\alpha\)-module via the adjoint
    representation.

    First consider
    \[
        M_\alpha
        :=
        \bigoplus_{c\in F} L_{c\alpha},
    \]
    where \(L_0=H\), and where \(L_{c\alpha}=0\) if \(c\alpha\notin\Delta\cup\{0\}\).
    This is an \(\mathfrak{s}_\alpha\)-submodule of \(L\). If
    \(x\in L_{c\alpha}\), then
    \[
        [h_\alpha,x]
        =
        c\alpha(h_\alpha)x
        =
        2c\,x.
    \]
    Hence \(L_{c\alpha}\) is the \(h_\alpha\)-weight space of weight \(2c\).
    By the finite-dimensional representation theory of \(\mathfrak{sl}_2\), all
    weights are integers. Therefore, whenever \(L_{c\alpha}\neq 0\), we have
    \[
        2c\in\mathbb{Z}.
    \]

    Now consider the submodule
    \[
        M_{\mathrm{even}}
        :=
        \bigoplus_{i\in\mathbb{Z}} L_{i\alpha}.
    \]
    All its \(h_\alpha\)-weights are even. Since every irreducible
    \(\mathfrak{sl}_2\)-module with even weights contains weight \(0\), each
    irreducible submodule \(S\) of \(M_{\mathrm{even}}\) satisfies
    \[
        S\cap H\neq 0.
    \]
    Choose \(0\neq h\in S\cap H\).

    If \(h\in\operatorname{Ker}\alpha\), then
    \[
        [x_\alpha,h]=[y_\alpha,h]=[h_\alpha,h]=0,
    \]
    so \(Fh\) is a trivial \(\mathfrak{s}_\alpha\)-submodule. Since \(S\) is
    irreducible, it follows that
    \[
        S=Fh.
    \]

    If \(h\notin\operatorname{Ker}\alpha\), then
    \[
        [x_\alpha,h]
        =
        -\alpha(h)x_\alpha\neq 0,
        \qquad
        [y_\alpha,h]
        =
        \alpha(h)y_\alpha\neq 0.
    \]
    Thus
    \[
        x_\alpha,y_\alpha\in S.
    \]
    Consequently
    \[
        h_\alpha=[x_\alpha,y_\alpha]\in S,
    \]
    and hence
    \[
        \mathfrak{s}_\alpha\subseteq S.
    \]
    Since \(\mathfrak{s}_\alpha\) is the irreducible adjoint module for itself,
    we get
    \[
        S=\mathfrak{s}_\alpha.
    \]
    Therefore
    \[
        M_{\mathrm{even}}
        =
        H+\mathfrak{s}_\alpha.
    \]
    It follows in particular that
    \[
        2\alpha\notin\Delta
        \qquad\text{and}\qquad
        \dim L_\alpha=1.
    \]
    Applying this to the root \(\frac{1}{2}\alpha\), if it existed, would imply
    \(\alpha\notin\Delta\), a contradiction. Hence
    \[
        \frac{1}{2}\alpha\notin\Delta.
    \]

    Next consider the submodule
    \[
        M_{\mathrm{odd}}
        :=
        \bigoplus_{i\in\mathbb{Z}} L_{\left(i+\frac12\right)\alpha}.
    \]
    All its \(h_\alpha\)-weights are odd. If this module were nonzero, then some
    irreducible summand would have odd weights, hence would contain weight \(1\).
    This would force
    \[
        L_{\frac12\alpha}\neq 0,
    \]
    contradicting the previous paragraph. Therefore
    \[
        M_{\mathrm{odd}}=0.
    \]
    Combining this with \(2c\in\mathbb{Z}\), we obtain
    \[
        F\alpha\cap\Delta=\{-\alpha,\alpha\},
        \qquad
        \dim L_\alpha=1.
    \]

    Now let \(\alpha,\beta\in\Delta\) with \(\beta\neq \pm\alpha\). Consider
    \[
        M_{\alpha,\beta}
        :=
        \bigoplus_{i\in\mathbb{Z}} L_{\beta+i\alpha}.
    \]
    This is again an \(\mathfrak{s}_\alpha\)-submodule of \(L\). If
    \(x\in L_{\beta+i\alpha}\), then
    \[
        [h_\alpha,x]
        =
        (\beta+i\alpha)(h_\alpha)x.
    \]
    Using
    \[
        \alpha(h_\alpha)=2
        \qquad\text{and}\qquad
        \beta(h_\alpha)=\frac{2(\alpha,\beta)}{(\alpha,\alpha)},
    \]
    we get
    \[
        (\beta+i\alpha)(h_\alpha)
        =
        \frac{2(\alpha,\beta)}{(\alpha,\alpha)}+2i.
    \]
    Since the weights of a finite-dimensional \(\mathfrak{sl}_2\)-module are
    integers, it follows that
    \[
        \frac{2(\alpha,\beta)}{(\alpha,\alpha)}\in\mathbb{Z}.
    \]

    Moreover, all these weights have the same parity. By the first part, each
    nonzero root space \(L_{\beta+i\alpha}\) is one-dimensional. Hence
    \(M_{\alpha,\beta}\) is irreducible as an \(\mathfrak{s}_\alpha\)-module.
    Therefore the weights occurring in \(M_{\alpha,\beta}\) form one complete
    \(\mathfrak{sl}_2\)-string. Thus there exist \(p,q\in\mathbb{Z}_{\geq 0}\)
    such that the roots of the form \(\beta+i\alpha\) are exactly
    \[
        \beta-q\alpha,\,
        \beta-(q-1)\alpha,\,
        \ldots,\,
        \beta-\alpha,\,
        \beta,\,
        \beta+\alpha,\,
        \ldots,\,
        \beta+p\alpha.
    \]
    The highest weight is
    \[
        \beta(h_\alpha)+2p,
    \]
    and the lowest weight is
    \[
        \beta(h_\alpha)-2q.
    \]
    Since the weights of an irreducible \(\mathfrak{sl}_2\)-module are symmetric
    about \(0\), we have
    \[
        \beta(h_\alpha)-2q
        =
        -\bigl(\beta(h_\alpha)+2p\bigr).
    \]
    Hence
    \[
        q-p
        =
        \beta(h_\alpha)
        =
        \frac{2(\alpha,\beta)}{(\alpha,\alpha)}.
    \]

    Finally, if \(\alpha+\beta\in\Delta\), then \(L_{\alpha+\beta}\neq 0\).
    In the irreducible \(\mathfrak{s}_\alpha\)-module \(M_{\alpha,\beta}\), the
    raising operator \(x_\alpha\) maps
    \[
        L_\beta
        \longrightarrow
        L_{\beta+\alpha}.
    \]
    Since \(\beta+\alpha\) is a weight, this map is nonzero. As
    \(L_{\beta+\alpha}\) is one-dimensional, we obtain
    \[
        [L_\alpha,L_\beta]
        =
        L_{\alpha+\beta}.
    \]
    This proves the proposition.
\end{proof}

Let \(\Delta\) be a root system in a Euclidean space \(E\). For
\(\alpha\in\Delta\), define the reflection
\[
    s_\alpha:E\longrightarrow E
\]
by
\[
    s_\alpha(x)
    =
    x-\frac{2(x,\alpha)}{(\alpha,\alpha)}\alpha.
\]
Then
\[
    s_\alpha(\alpha)
    =
    \alpha-\frac{2(\alpha,\alpha)}{(\alpha,\alpha)}\alpha
    =
    -\alpha.
\]
Moreover, if
\[
    H_\alpha
    :=
    \{\,x\in E\mid (x,\alpha)=0\,\},
\]
then
\[
    s_\alpha(x)=x
    \qquad
    \text{for all }x\in H_\alpha.
\]
Thus \(s_\alpha\) is the reflection with respect to the hyperplane
\(H_\alpha\), which has codimension \(1\) in \(E\).

\begin{lemma}
    For every \(\alpha\in\Delta\), one has
    \[
        s_\alpha(\Delta)=\Delta.
    \]
\end{lemma}

\begin{proof}
    Since \(s_\alpha^2=\operatorname{id}_E\), the map \(s_\alpha\) is a bijection.
    Hence it suffices to prove that
    \[
        s_\alpha(\Delta)\subseteq \Delta.
    \]

    Let \(\beta\in\Delta\). By the root-string property, there exist
    \(p,q\in\mathbb{Z}_{\geq 0}\) such that the roots of the form
    \(\beta+i\alpha\) are precisely
    \[
        \beta-q\alpha,\,
        \beta-(q-1)\alpha,\,
        \ldots,\,
        \beta,\,
        \ldots,\,
        \beta+p\alpha,
    \]
    and
    \[
        q-p
        =
        \frac{2(\beta,\alpha)}{(\alpha,\alpha)}.
    \]
    Therefore
    \[
        s_\alpha(\beta)
        =
        \beta-\frac{2(\beta,\alpha)}{(\alpha,\alpha)}\alpha
        =
        \beta-(q-p)\alpha
        =
        \beta+(p-q)\alpha.
    \]
    Since
    \[
        -q\leq p-q\leq p,
    \]
    the root-string property implies that
    \[
        \beta+(p-q)\alpha\in\Delta.
    \]
    Hence \(s_\alpha(\beta)\in\Delta\). Since \(\beta\in\Delta\) was arbitrary,
    we have
    \[
        s_\alpha(\Delta)\subseteq\Delta.
    \]
    As \(s_\alpha\) is bijective, it follows that
    \[
        s_\alpha(\Delta)=\Delta.
    \]
\end{proof}

